%======================================================================
\section{記号添字規約}
\label{app:notation}

本付録では、本論文の記号に付す添字の位置と順序を、
添字が担う役割の種類ごとに一つに定め、
同じ種類の添字はどの節でも同じ位置に置く。
本付録の規約は記法上の約束であり、
仮定一覧には含めない。
各記号の数学的内容は、それぞれの初出節の定義による。

%----------------------------------------------------------------------
\subsection*{記号添字規約 NIC（添字の位置・順序・書体）}
\label{conv:NIC}

\begin{description}

  \item[(N1) 下付きの識別列]
    下付きには識別列を置く。
    対象を識別する添字は、すべて下付きの識別列に置き、
    \[
      \text{種類}\;\to\;\text{個体}\;\to\;\text{成分}\;\to\;\text{窓}\;\to\;\text{時刻}
    \]
    の順に並べる。
    種類の位置には、説明的ラベル、
    ラベル集合上を走る種類変数、
    および写像・評価関数の依存先を示すパラメータを置き、
    二つ以上が並ぶ場合は、量の直接の種別を先に、
    水準等の修飾を後に置く。
    個体の位置には、ACS個体および外部対象の識別子と、その順序対を置く。
    成分の位置には、TIPAの成分番号、状態座標の番号および組の成分の番号を置く。
    窓の位置には観測窓の長さを置き、
    時刻は下付きの列の最後に置く。
    固定の説明的ラベルを、下付きの種類の位置に置くか、
    (N2) の上付きの変種ラベルとするかは、
    その記号の初出の定義で定め、以後変えない。
    下付きの種類の位置に置いたラベルは識別列の一部であり、
    個体・成分・窓・時刻の添字はその後に続く。
    例えば $d_{\mathrm{eff},i,t}$ は個体 $i$ の時刻 $t$ における
    $d_{\mathrm{eff}}$ を、
    $\bar Q_{i\leftarrow j,W_Q,t}$ は順序対 $i\leftarrow j$ についての
    窓幅 $W_Q$ の時刻 $t$ における値を、
    $E_{\mathrm{ext},\mathrm{org}}$ は組織水準の外部環境状態空間を表す。

  \item[(N2) 上付き]
    上付きに置くのは、次のいずれかに限る。
    \begin{enumerate}[label=(\alph*)]
      \item 演算：冪、逆写像、補集合、参照点、転置、関数空間の指数、
      および括弧付きで表す反復の段数（$\mathrm{Crit}^{(1)}$、$\mathrm{Viab}_f^{(\nu)}(K)$ 等）。
      括弧付きの上付きを段数以外の意味で用いる場合は、使用箇所で宣言する。
      \item 変種ラベル：同じ基本記号で表される対象の変種を区別する、
      固定した説明的ラベル（$E^{\mathrm{bg}}$、$U^{\mathrm{rob}}$、$Q^{\mathrm{inf}}$ 等）。
      変種ラベルは集合上を走る変数ではなく、
      識別列に参与しない。
      変種ラベルを付けた記号には、(N1) の識別列を下付きに続けて置く
      （$Q^{\mathrm{inf}}_{i\leftarrow j,t}$ 等）。
    \end{enumerate}
    個体・成分・窓・時刻の識別は上付きに置かない。
    総和、積、合併等の大型演算子の範囲指定は通常の位置に置き、
    本項の対象としない。

  \item[(N3) 括弧]
    括弧内には、数学的対象を入力とする写像、対応または演算の入力を置く。
    括弧を用いるのは、値がその入力によって定まる場合に限る。
    入力は、宣言された集合の元であるか、
    または (N9) の構造型を満たす任意の構造に対して定まる演算における
    構造そのものである。
    例えば $f(x,u,e)$、$A(x,e)$、$E_{\mathrm{req}}(x)$、$\mathrm{Viab}_f(K)$、
    $S_{\mathrm{struct}}(\Sigma_i)$ の括弧内はこの意味での入力である。

  \item[(N4) 識別と入力]
    値が、識別された個体の実現、すなわちその個体の状態列や環境実現に依存する量は、
    構造を入力とする関数形では書かず、
    識別子と時刻を下付きに置く。
    このため、写像を定義したうえで、その時刻ごとの評価値を
    \[
      A_{i,t}:=A_i(x_{i,t},e_{i,t})
    \]
    のように下付きで表す。
    時刻評価値の記号は各節で定義を与えたものに限り、
    それ以外は $E_{\mathrm{req}}(x_t)$ のように
    写像に状態を入力した形で書く。

  \item[(N5) 個体と水準]
    個体ACSの識別子 $i,i'$ は、\S\ref{sec:org_acs} の識別子集合 $\mathbb I$ の元とする。
    外部対象は $Y_j\in\mathcal Y_{-i}$ と書き、
    添字 $j$ は外部対象を区別する識別子とする。
    $-i$ は $i$ を除く他者を表す派生記法であり、
    $-i$ 自身は識別子ではない。
    個体 $i$ の構造を $\Sigma_i$、
    その成分を $X_i,U_i,E_i,f_i,\Pi_i,M_i,\mathrm{Update}_i$、
    時刻 $t$ の状態を $x_{i,t}$、$M_{i,t}$、
    個体 $i$ に相対的な外部環境状態を
    $e_{i,t}\in E_{\mathrm{ext},i}$ と書く。
    組織水準等の水準名は種類のラベルとして扱い、
    $\Sigma_{\mathrm{org}}$ と書く。
    具体例では、宣言したうえで個体を数字で識別してよい。
    複数個体を扱う文脈では、個体の位置を省略しない。

  \item[(N6) 時刻]
    本論文では、時刻は列または過程の添字として下付きの最後に置き、
    写像の入力としては扱わない。

  \item[(N7) 書体]
    値が走る添字は斜体で書き、
    説明的ラベル、変種ラベルおよび構造部位名は立体で書く。
    ただし、他の記号を名指すラベル（$\mathbb T^{Q}$ の $Q$ 等）は、
    名指す記号の書体に従う。
    空間および集合は大文字で、その元および状態は小文字で書くことを原則とし、
    内部モデルの状態 $M_t$ はその例外とする（(N9)）。
    上付きと下付きの区別には、
    テンソル解析における反変・共変のような
    変換上の意味をもたせない。

  \item[(N8) 省略と式の成分]
    単一個体のみを扱う文脈では個体の添字を省略する。
    下付きが五つ以上並ぶ場合は、
    固定する添字を段落の冒頭で宣言したうえで、
    その添字を省略する。
    式の成分は $[f(x,u,e)]_k$ のように角括弧を用いて表す。

  \item[(N9) 構造と状態]
    定義~\ref{def:acs} が定める7タプルの構造型を抽象構造型と呼び、
    個体の構造 $\Sigma$、$\Sigma_i$ をその実現とする。
    抽象構造型を集合として扱うことは要求せず、
    個体 $\Sigma_i$ を抽象構造型の元とは書かない。
    内部モデルについては、
    構造 $M$（第 $c$ 機構 $M_c$）、
    状態空間 $\mathcal M$（成分 $\mathcal M_c$）、
    状態 $M_t\in\mathcal M$（成分 $M_{c,t}\in\mathcal M_c$）を区別する。
    構造に時刻添字を付けた $\Sigma_{i,t}$ は構造の時間変化を記述するための予約記号とし、
    本論文では用いない（\S\ref{subsec:limitations_list} の L16）。

  \item[(N10) 文字の予約]
    添字に用いる文字は、種類ごとに次の表のとおり予約する。
    予約は下付きの添字として用いる文字に適用する。
    集合の元を表す束縛変数を使用箇所で宣言して用いる場合
    （$q\in\mathrm{Req}_t$ 等）は、本項の対象としない。

\end{description}

\begin{center}
\begin{tabular}{ll}
  \hline
  文字 & 種類 \\
  \hline
  $i,\ i'$ & 個体ACSの識別子（$\mathbb I$ の元） \\
  $j$ & 外部対象の識別子 \\
  $c$ & TIPAの成分番号 \\
  $k$ & 状態座標の番号および組の成分の番号 \\
  $a$ & 行動応答成分の枠（$\mathcal A_O$ の元） \\
  $\kappa$ & 識別の行のラベル（$\mathcal R_\xi$ の元） \\
  $\iota$ & 偏向性の添字（$\mathcal J_\beta$ の元） \\
  $\ell$ & 機能的記憶の種類 \\
  $q$ & 健全度の種類 \\
  $T,\ L$ & 観測窓の長さ \\
  $t,\ s,\ r$ & 時刻 \\
  $\nu$ & 段数（上付きの括弧内） \\
  \hline
\end{tabular}
\end{center}

数字の下付きは、番号で区別して定義した特定の対象
（$M_1$、$P_2$、$\mathsf Z_3$、$t_0$、$\varepsilon_0$ 等）、
または (N5) により宣言した具体例の個体を表す。
軌道 $\sigma$ は軌道記録空間の元として括弧内に置く。
$n$、$m$ は自然数を表す記号として使用箇所で意味を宣言し、
添字には用いない（$\prod_{i=1}^{n}$ のような範囲指定には用いる）。

\begin{description}
  \item[(N11) 書体の使用規則]
    ひげ文字（フラクトゥール体）は、
    原則として母集合・台集合および集合族・写像族などの集合的対象に用いる
    （$\mathfrak K$、$\mathfrak P$ 等）。
    個々の写像、演算子および評価規則には、
    名称型（$\mathrm{rd}$、$\mathrm{Rule}$ 等）または通常の数式記号を用いる。
    書体の違いは記述対象を区別するための表記上の規約であり、
    書体のみから追加的な数学的性質を導かない。
\end{description}
